\section{Resultados, Validacion y Evidencias}

\subsection{Resultados tecnicos alcanzados}
Al cierre del periodo de practicas se consideran alcanzados los siguientes resultados:
\begin{itemize}[leftmargin=1.5em]
    \item Estabilizacion inicial del comportamiento en VR tras ajustes de configuracion y escena.
    \item Sistema de comportamientos orientado a crecimiento mediante estructura de estados/transiciones mas clara.
    \item Conexion web reforzada con comprobaciones explicitas para eventos desconocidos o no implementados.
    \item Cobertura funcional de conflictos existente verificada e incremento del catalogo con cinco nuevos conflictos (TEA/TDAH).
    \item Integracion de animaciones y prefabs para representar los nuevos casos en escena.
    \item Ejecucion validada en Meta Quest 3 sin bloqueos graves observados en el flujo principal.
    \item Mejora operativa para el equipo gracias a seleccion de servidores de prueba desde interfaz.
\end{itemize}

\subsection{Criterios de verificacion aplicados}
La validacion no se realizo como una prueba unica al final, sino mediante controles recurrentes durante cada iteracion. Los criterios principales fueron:
\begin{itemize}[leftmargin=1.5em]
    \item Correctitud funcional: el conflicto solicitado debe derivar en la transicion y animacion esperadas.
    \item Robustez: ante mensajes anomalos o incompletos, el sistema debe registrar el error y evitar estado inconsistente.
    \item Reproducibilidad: los casos de prueba deben poder repetirse entre editor y build.
    \item Integracion: cambios de codigo, prefabs y animacion deben mantenerse sincronizados.
    \item Compatibilidad objetivo: flujo estable en dispositivo Quest 3.
\end{itemize}

\subsection{Impacto en el flujo de equipo}
Las mejoras introducidas han tenido impacto directo en la colaboracion:
\begin{itemize}[leftmargin=1.5em]
    \item Menor coste de depuracion por disponer de logs y aserciones mas explicitas.
    \item Mayor rapidez en pruebas cruzadas al poder alternar entorno de servidor.
    \item Menor friccion para incorporar nuevos conflictos al existir una ruta de integracion mas definida.
\end{itemize}

\subsection{Limitaciones detectadas y trabajo futuro}
Aunque los objetivos principales se completaron, quedan lineas de continuidad recomendadas:
\begin{itemize}[leftmargin=1.5em]
    \item Formalizar una bateria de pruebas automatizadas de regresion para conflictos y transiciones.
    \item Profundizar en metricas cuantitativas de rendimiento en Quest 3 (FPS medio, p95/p99 de frametime, consumo de memoria).
    \item Consolidar documentacion funcional de los nuevos conflictos para perfiles no tecnicos del equipo.
    \item Evaluar herramientas de telemetria para diagnostico remoto de sesiones.
\end{itemize}

\subsection{Resumen ejecutivo de hitos}
\begin{center}
\begin{tabular}{p{0.26\linewidth} p{0.67\linewidth}}
\toprule
\textbf{Hito} & \textbf{Resultado} \\
\midrule
Optimizar base VR & Reduccion de fricciones visuales y ajustes de entrada/rastreo para estabilidad de uso. \\
Dise\~nar base de comportamientos & Estructura de estados/transiciones preparada para ampliar casos no normativos. \\
Reforzar conexion web & Mayor control de errores, deteccion de eventos desconocidos y trazabilidad de fallos. \\
Extender conflictos & Verificados 3 existentes y a\~nadidos 5 nuevos (TEA/TDAH) con soporte de pipeline. \\
Integrar contenido & Animaciones y prefabs adaptados para representar correctamente nuevos escenarios. \\
Exportar a Quest 3 & Build funcional validada en dispositivo objetivo. \\
\bottomrule
\end{tabular}
\end{center}